% Program Map — shared front-matter page for all four volumes.
% Canonical source: style/program_map.tex; copied into each volume by `make sync-style`.
\chapter*{The Program Map}
\addcontentsline{toc}{chapter}{The Program Map}

This volume is one of four. Together they cover the surfaces a Staff-level machine-learning loop tests --- breadth, depth, coding, system design, and the staff signals themselves --- without repeating each other. Each topic has exactly one \emph{canonical home}: the volume that carries the derivation and the deep treatment. Everywhere else, that topic appears as a short summary with a pointer. If a chapter seems to stop early on something, it is because the full treatment lives elsewhere and the pointer will say where.

\section*{The Four Volumes}

\begin{tabularx}{\textwidth}{@{}l l X@{}}
\toprule
\textbf{Vol.} & \textbf{Title} & \textbf{What it owns} \\
\midrule
I & Deep Learning Essentials & Mathematical and learning-theory foundations, losses, architectures, training and optimization at scale, pretraining, GPU performance, inference and serving, agents, safety, evaluation, production, and the interview-craft chapters (coding rounds, the staff loop). \\
II & NLP Essentials & Classical and statistical NLP, word representations, sequence models, transformer derivations, pretraining and adaptation, in-context learning and test-time compute, alignment practice, RAG, the benchmark landscape, production NLP, safety and ethics. \\
III & Search \& Recommendation Essentials & Query understanding, lexical and vector retrieval, neural retrieval and reranking, learning to rank, recommender systems, ranking metrics and experimentation, production retrieval and RAG serving. \\
IV & Conventional ML Essentials & Linear models, trees and ensembles (the GBDT home), unsupervised methods, probabilistic foundations, applied craft, experimentation and causal inference, time series, and from-scratch classical implementations. \\
\bottomrule
\end{tabularx}

\section*{Cross-Volume References}

PDFs cannot hyperlink each other, so cross-volume pointers are written in brackets with the volume and chapter number: \textbf{[DL 5]} is Volume I, Chapter 5; \textbf{[NLP 9]} is Volume II, Chapter 9; \textbf{[SR 4]} is Volume III, Chapter 4; \textbf{[CML 2]} is Volume IV, Chapter 2. References \emph{within} a volume use ordinary numbered cross-references.

\section*{Where Each Overlapping Topic Lives}

These are the zones where two volumes could plausibly claim a topic. The table is the contract:

\begin{tabularx}{\textwidth}{@{}l l X@{}}
\toprule
\textbf{Topic} & \textbf{Canonical home} & \textbf{What the other volumes keep} \\
\midrule
Transformer derivations & [NLP 5] & [DL 5] keeps the quantitative systems view: parameter and FLOP counting, KV-cache byte math, MLA, FlashAttention, sparse attention. \\
Embedding theory & [NLP 3] & [DL 6] keeps production diagnostics: collapse, spectral checks, tokenization interaction, memory arithmetic. \\
Tokenization & [NLP 12] & [DL 9] keeps interaction notes; [DL 28] keeps the BPE implementation drill. \\
Metric-learning training & [DL 7] & [SR 4] keeps retrieval-system training recipes and serving. \\
ANN and vector indexes & [SR 3] & [DL 7] keeps a one-page summary and points here. \\
Retrieval system design & [SR 4] & One canonical design question; [DL 6] and [NLP 3] point to it. \\
Ranking metrics (NDCG, MRR, MAP) & [SR 7] & [DL 23] keeps one-line catalog entries. \\
Recommender systems & [SR 6] & [DL 12] is a short bridge chapter. \\
RAG pipeline & [NLP 9] & [SR 8] owns the retrieval stage; [DL 21] owns the context-vs-retrieval decision. \\
Long context & [DL 21] & [NLP 5] keeps positional-encoding schemes. \\
PEFT mechanics & [DL 16] & [NLP 6] keeps the adaptation-lifecycle narrative. \\
Pretraining data and run recipes & [DL 25] & [NLP 1--2] keep the classical instruments that run inside those pipelines. \\
RL foundations and reasoning & [DL 17] & [NLP 8] owns alignment practice and RLVR recipes. \\
Alignment pipeline & [NLP 8] & [DL 17] keeps the RL mathematics. \\
Agent architectures and evaluation & [DL 27] & [NLP 7] owns prompting and tool-calling formats. \\
Technical safety methods & [DL 20] & [NLP 13] keeps NLP-specific harms and evaluation. \\
Hallucination and faithfulness & [NLP 13] & [DL 20] keeps the safety framing for deployed and agentic settings; [NLP 9] owns RAG grounding. \\
Inference and serving & [DL 19] & [NLP 12] keeps the LLM application layer. \\
LLM evaluation statistics & [DL 23] & [NLP 10] owns the benchmark landscape and judge practice. \\
GBDT internals & [CML 2] & [SR 5] keeps LambdaMART's ranking-specific view. \\
A/B and experimentation statistics & [CML 6] & [SR 7] keeps ranking-specific experimentation (interleaving). \\
\bottomrule
\end{tabularx}

\section*{Reading Paths}

\textbf{Breadth round.} [DL 1--4] foundations, [DL 23] metrics, [DL 24] decision frameworks, the whole of Volume IV, [NLP 1--3], and [SR 1, 6, 7].

\textbf{LLM / frontier-lab depth.} [NLP 5] derivations, [DL 5] quantitative systems, [NLP 6] pretraining and adaptation, [DL 25] pretraining at scale, [DL 26] GPU performance, [DL 15] optimization, [DL 14] efficient architectures, [NLP 7] in-context learning and test-time compute, [DL 17] and [NLP 8] for RL and alignment, [DL 19] inference, [DL 21] long context, [DL 27] agents, [DL 20] and [NLP 13] safety.

\textbf{Search and recommendation depth.} [CML 2] and [CML 6] first, then Volume III cover to cover, then [DL 6--8] for embeddings and contrastive training, [NLP 9] for the RAG pipeline, [DL 21] for the context-versus-retrieval decision, and [DL 19] for serving.

\textbf{Coding round.} [DL 28] for strategy and the worked canon, with the runnable drills in \texttt{drills/}.

\textbf{The loop itself.} [DL 29] for level calibration, project narratives, behavioral scenarios, and research taste.
